%% ============================================================
%% Chapter 20: Rebel Without a Cost
%% ============================================================

\chapter{Rebel Without a Cost}
\label{ch:ch20-rebel-without-a-cost}
\section{The Retrospective Problem}

An earlier internal feasibility study, developed under the working name Ayian, proposed converting recorded contribution into a work-backed, persona-bound cryptocurrency, treating recording, recognizing, rewarding, and authorizing as stages of one continuous mechanism. It compared Layer-1 and Layer-2 chains, separated fungible rewards from non-fungible personas, proposed token-bound accounts, and attempted to reduce whale power through a hybrid voting formula. Its governing aspiration was valuable: useful work should count, durable contribution should not disappear, and those who sustain a system should have a meaningful voice in its continuation. None of what follows disputes that aspiration. It disputes the mechanism chosen to realize it.

The difficulty lies not primarily in the selected chain, token standard, or voting exponent. It lies in the compression of several different judgments into a single scalar object. The proposed system turns heterogeneous contributions into a contribution score, converts that score into token issuance, stores the tokens under a persistent persona, and then feeds both stake and score back into governance. Documentary history becomes economic value; economic value becomes political power; identity becomes the container binding them together. This is a specific instance of the same log-admissibility-without-positive-productivity confusion Chapter~\ref{ch:ch19-ledgers-without-value} formalized: a contribution ledger can be perfectly consistent, every record accurate, every score computed as specified, while the conversion from record to price to authority remains structurally unwarranted at every step.

This is an attractive pipeline because each step appears, locally, to preserve the legitimacy of the step before it. Recording seems to warrant recognizing; recognizing seems to warrant rewarding; rewarding seems to warrant a say in what happens next. The appeal is not naive. It is the same appeal that makes seniority, credentialism, and inherited reputation feel natural in institutions that never touched a blockchain. What this chapter adds is an account of exactly where, and exactly why, each of those local inferences fails, stated precisely enough to be checked rather than merely felt.

\section{Six Acts That Are Not One Act}

The central formal claim is a chain of non-implications, running from the fact that an event occurred to the much stronger claim that it should govern what happens next:
\[
\mathrm{Record}(x) \not\Rightarrow \mathrm{Verify}(x) \not\Rightarrow \mathrm{Recognize}(x) \not\Rightarrow \mathrm{Reward}(x) \not\Rightarrow \mathrm{Authorize}(x) \not\Rightarrow \mathrm{Continue}(x).
\]
Verification is placed here deliberately, and separately from recognition, because the two do different jobs. Verification asks whether an event actually occurred as described, a factual, evidentiary question. Recognition asks whether the event, having occurred, is regarded as valuable, a normative, social question. Conflating them lets ``recognized'' silently mean both ``confirmed to have happened'' and ``endorsed as good,'' which obscures the case that matters most: an act can be accurately recorded and independently verified without being recognized as valuable at all, because verification establishes only that something happened, not that it was welcome.

A contribution may therefore be recorded without being verified. It may be verified without being recognized as valuable, a harmful or mistaken act can be confirmed to have occurred and still deserve no endorsement. It may be recognized without being rewarded, gratitude is not payment. It may be rewarded without conferring permanent standing, payment for a job does not establish general expertise. It may establish standing in one domain without granting authority in another. And it may have been an admissible basis for a decision yesterday without remaining one today. Each arrow in the chain corresponds to an institution people already build to keep the two sides separate, precisely because collapsing them causes recognizable harm: archives distinguish preservation from endorsement, fact-checking distinguishes verification from approval, mutual aid networks distinguish gratitude from payment, professional licensing distinguishes payment from standing, federalism distinguishes domain expertise from general authority, and constitutional term limits distinguish past merit from future permission. A system that mints one convertible token across all six acts does not innovate past these institutions. It discards the distinctions they were built to protect. This six-object chain is the contribution-pipeline instance of the same $H \not\Rightarrow D \not\Rightarrow R \not\Rightarrow C$ structure Chapter~\ref{ch:ch10-clio-architecture} developed for agentic systems and Chapter~\ref{ch:ch11-constitutional-space-of-refusal} restated at the constitutional layer: verification constrains continuation, but does not authorize commitment, whether the acting party is an AI agent or a community of contributors.

The central revision is therefore architectural: Ayian should not begin with the question, ``How do we mint value from work?'' It should begin with the more conservative question, ``How do we prevent work, provenance, refusal, and obligation from being erased?''

\section{Why No Universal Scalar Is Canonical}

Work is heterogeneous: repairing a bridge, moderating a dispute, discovering a theorem, and caring for a sick person do not occupy a common natural scale, and any scalar conversion is imposed rather than discovered. This can be stated more precisely than that loose intuition, strong enough to carry the chapter's actual claim: a contribution record with genuinely incomparable elements does not, by itself, pick out one correct total ranking.

\begin{definition}[Contribution preorder and quotient order]
Let $C$ be a set of contribution records. For a declared context $\kappa$, specifying a domain, an affected population, an evidentiary standard, and a set of relevant dimensions, write $c_1 \preceq_\kappa c_2$ when every comparison admitted by $\kappa$ judges $c_2$ at least as well supported as $c_1$, with no relevant dimension on which $c_1$ is preferred. Contributions equivalent in both directions are identified, $c_1 \sim_\kappa c_2 \iff c_1 \preceq_\kappa c_2 \land c_2 \preceq_\kappa c_1$, and the quotient $C/{\sim_\kappa}$ is then partially ordered. Two records with no dominance relation in either direction remain incomparable in this quotient order.
\end{definition}

A scalar valuation $v : C \to \mathbb{R}$ is isotone with respect to $\preceq_\kappa$ when $c_1 \preceq_\kappa c_2$ implies $v(c_1) \leq v(c_2)$. Isotonicity alone is weak: an isotone $v$ may assign equal values to incomparable contributions, which is a defensible and honest choice, since it registers that the record does not distinguish them. The stronger and more common design choice, made implicitly by any system that produces a single ranked leaderboard, requires $v$ to be strict: to assign distinct values to any two contributions that are not identified under $\sim_\kappa$. A strict isotone valuation induces a total order on its image, precisely a linear extension of the quotient order.

\begin{proposition}[Non-canonicity of total contribution rankings]
Let $(C/{\sim_\kappa}, \preceq_\kappa)$ contain an incomparable pair $c_1, c_2$. Then $\preceq_\kappa$ does not determine a unique total extension of that pair: there exists a linear extension in which $c_1$ precedes $c_2$, and another, equally consistent with $\preceq_\kappa$, in which $c_2$ precedes $c_1$. Consequently, any scalar system required to rank every pair of contributions must either preserve some incomparabilities as ties, or introduce a comparison that $\preceq_\kappa$ alone does not warrant.
\end{proposition}

Let $c_1$ and $c_2$ be incomparable under $\preceq_\kappa$. Adjoining $c_1 \prec c_2$ and taking transitive closure introduces no cycle, since a cycle would require $c_2 \preceq_\kappa c_1$ already, contradicting incomparability; the order-extension theorem then guarantees a linear extension realizing $c_1 \prec c_2$. Reversing the construction produces, by the identical argument, a linear extension realizing $c_2 \prec c_1$. Both preserve every comparison already present in $\preceq_\kappa$, so neither is disqualified by the original record, and the record itself supplies no fact that selects between them. This establishes non-canonicity, not impossibility: a total ranking can certainly be constructed, and in a finite contribution set there are only finitely many linear extensions to choose among. The point is narrower and sufficient: no fact contained in the contribution record itself selects which of those finitely many extensions is the right one. A scalar system that ties incomparable contributions is being honest about this; one that breaks the tie with a specific numeric ranking has quietly imported an institutional judgment and presented it as a discovery about the contributions themselves.

The severity of this non-canonicity is governed by the order's width, the maximum size of an antichain, a set of contributions no two of which are comparable under $\preceq_\kappa$. If a contribution poset contains an antichain of size $w$, its members are compatible with all $w!$ relative orderings. This motivates a design choice: rather than storing a single number that discards the antichain structure the poset actually has, preserve a structured record $r(c) = \langle a, t, d, e, p, o, q, s \rangle$, where $a$ is the actor, $t$ the time interval, $d$ the domain, $e$ the evidence, $p$ the provenance, $o$ the affected objects or communities, $q$ the contestable claims made about the contribution, and $s$ its present status. Two contributions incomparable under $\preceq_\kappa$ because they differ in domain and affected population remain visibly, honestly incomparable in $r(c)$, rather than being forced through a projection that manufactures a rank.

The order-theoretic argument establishes that ranking adds unwarranted comparisons. A separate argument establishes that a transferable reward built on such a ranking cannot remain a faithful indicator of the responsibility it was meant to track.

\begin{proposition}[Transfer-fidelity conflict]
Suppose a system issues a transferable unit for contribution $c$, that possession of that unit increases governance authority, and that governance authority is required to track the responsibility or competence evidenced by $c$. These three requirements cannot be jointly guaranteed.
\end{proposition}

Let $a$ receive a unit $\tau$ for performing $c$, and let $b$ subsequently acquire $\tau$ without acquiring $a$'s responsibility for, or competence demonstrated by, $c$. Transferability permits this change of possession. Governance convertibility means possession of $\tau$ increases $b$'s authority regardless of how $\tau$ was acquired. Fidelity requires that authority track the relation to $c$ that justified $\tau$'s issuance, a relation $b$ does not hold. Hence $b$'s increased authority is not warranted by fidelity, and transferability, convertibility, and fidelity cannot all hold simultaneously. Together, these two propositions support the chapter's title theme directly: even setting aside whether a contribution score can be canonical, making that score's reward transferable severs the one relation, between the person acting and the authority the action is meant to justify, that a governance system built on ``proof of contribution'' was supposed to preserve.

\section{Prior Art: Ledgers That Already Refuse Conversion}

Systems that separate documentary record from price and from authority already exist, in some cases predating any blockchain by decades. Wikipedia's page-history mechanism normally retains and exposes every edit alongside a separate and much smaller mechanism, administrator status, that confers operational authority; edit count alone grants no vote and no formal reputation score, and even a very large edit history confers no automatic operational authority absent a distinct nomination process that can be, and sometimes is, refused. Git's commit history performs a structurally similar separation: content-addressing plus optional signing plus independent mirroring makes silent retroactive alteration conspicuous rather than impossible, but write access, merge rights, and governance votes are typically granted through a separate maintainer process rather than by raw commit count. Local Exchange Trading Systems and time banking keep a running account of hours or services exchanged without treating that account as a source of political authority over the community that issued it, arriving independently at a similar architectural instinct: record contribution, but constrain what the record is allowed to become.

Academic citation practice presents the clearest cautionary case in the other direction, showing what happens when the separation this chapter recommends is allowed to erode. A citation is, in origin, a record; citation-derived metrics such as the journal impact factor and the h-index are a further step removed, aggregations of that record into a score, and the metrics-reform literature has documented how such scores, once adopted for hiring and funding decisions, distort the incentives of the research they were meant merely to describe. A citation record and a citation-derived ranking are not the same object, and treating the second as an unproblematic summary of the first reproduces, in an academic register, the same collapse this chapter argues against in a token-governed one.

\section{From a Value Ledger to a Documentary Ledger}

The proposed ledger is append-only but not declaratively omniscient. It records claims and their provenance rather than asserting one final account of worth. A contribution event may be represented as
\[
e_i = \langle \mathrm{id}, \mathrm{actor}, \mathrm{action}, \mathrm{object}, \mathrm{time}, \mathrm{evidence}, \mathrm{attestations}, \mathrm{contestations}, \mathrm{scope} \rangle.
\]
The history after $n$ events is $H_n = H_0 \parallel e_1 \parallel \cdots \parallel e_n$, where append-only retention means that later judgments do not silently rewrite earlier evidence. This needs qualification, since an unconditional prohibition on removing any byte conflicts with legitimate privacy, safety, and legal obligations: a later judgment may supersede, contest, seal, or restrict access to an earlier record, but it does not silently replace that record, and where privacy, safety, consent, or law requires underlying material to be sealed or removed, the ledger retains, where permissible, a minimal audit event recording that a governed removal occurred, by whose authority, under which rule, and with what scope, without necessarily retaining the removed content itself.

Operative state is separate from the retained history, $\sigma_n = J(H_n, \kappa_n)$, where $J$ is a revisable judgment procedure and $\kappa_n$ is the current context, exactly the $\sigma$-versus-$H$ separation Chapter~\ref{ch:ch10-clio-architecture} formalized as the central discipline of the CLIO architecture. The same history can therefore support different admissible continuations in different domains without contradiction: a medical contribution need not confer software-governance authority; an early architectural contribution need not grant perpetual control; a rejected proposal can remain historically important without becoming operative.

A transferable unit of value can itself be understood as a documentary record from which nearly every field has been removed. A dollar does not ordinarily retain who performed the underlying work, who was affected by it, what obligations were discharged in exchange for it, what harms accompanied its production, or whether its present holder bears any relation at all to the event that caused the unit to exist. Money is not the opposite of a ledger. It is an amputated one. Fungibility is achieved by deleting provenance until units produced by radically different histories become operationally interchangeable. A ledger without value reverses that compression: it preserves reasons while declining to manufacture a universal, frictionless permission from them, making it less liquid but more articulate. The earlier Ayian architecture, in minting a new transferable token from contribution records, was choosing amnesia while believing it was choosing memory.

\section{Persona Without Ownership}

The original Ayian architecture placed an ERC-20 balance inside an account controlled through a Persona NFT. This elegantly binds assets to an identity object, but it also creates a dangerous ambiguity: if the persona token is transferable, then identity and its accumulated history can be sold; if it is non-transferable, then control recovery, coercion, succession, and erroneous binding become difficult governance problems the original study did not resolve.

A persona in a ledger without value is not property. It is a continuity relation among cryptographic control, documentary history, social recognition, and the continuing consent of the subject. No single token constitutes the person. Keys may rotate; attestations may expire; a person may withdraw from an institution while retaining the record of what occurred there. The persona layer should satisfy four constraints: control must be recoverable without selling identity; claims must be scoped and revocable by their issuers without erasing their historical issuance; the subject must be able to refuse new associations with their self-presented identity; and the ledger must distinguish identity continuity from institutional membership, since leaving a system is not the destruction of a person, and retaining history is not permission to keep governing them.

This last constraint requires a distinction the earlier draft elided: control recovery is not the same operation as succession. Recovery asserts continuity of the same subject, a new key continuing the same persona. Succession transfers a role, custody, or unresolved obligation from one subject to a different one entirely, and must not be mistaken for the predecessor's identity passing along with it: a successor inherits stewardship of what a persona left behind, not the standing to have been that persona. A recovery mechanism built only around ``prove you are still you'' is vulnerable to an adversary who has compromised the subject's other credentials, or to a subject coerced into authorizing a transfer they do not actually consent to; a recovery process with a mandatory delay window, notification to previously designated attestors, and a contestation period before any high-consequence action takes effect gives a coerced or compromised subject a real chance to intervene before the action becomes irreversible.

A subject cannot reasonably veto every adverse claim made about them, that would let anyone erase accountability simply by refusing to acknowledge it, but neither should an unverified allegation appear in the record with the same standing as an accepted fact. The architecture should therefore distinguish subject-approved associations, which the subject has affirmatively accepted; third-party attributed claims, which remain visible and attributable to their source but are marked as contested until resolved; and institutionally admitted findings, which have passed through some accountable adjudication process. Consent governs the first category. It does not, and should not, govern the second or third, which is what preserves the ledger's usefulness for accountability rather than converting it into a subject-curated profile.

\section{Governance After the Voting-Power Formula}

The corrected Ayian formula was $P_i = \sqrt{S_i}\, C_i$, where $S_i$ is staked Ayian and $C_i$ is lifetime contribution score. The square root improves the formula relative to linear stake by diminishing the marginal power of large holders. It does not, however, resolve the underlying constitutional problem. Multiplication means that economic stake and recorded contribution validate one another: a high value in either dimension magnifies the other, while a zero in either annihilates the person's voice, and it turns a lifetime record into a permanent political multiplier, allowing old contribution to harden into office. The term ``quadratic voting'' is also imprecise here; taking $\sqrt{S_i}$ applies a concave weighting to stake, while quadratic voting ordinarily concerns the quadratic cost of purchasing multiple votes. More importantly, if $C_i$ is obtained by reducing heterogeneous contribution records to a universal score, it is one particular resolution of the non-canonical choice the previous section's proposition describes, a strict isotone valuation that has broken every tie the underlying order left open, so multiplying by $\sqrt{S_i}$ scales economic stake by comparisons that were introduced during the scalarization of history, not supplied by the history itself.

A second, independent problem concerns how $C_i$ behaves over time. If lifetime contribution is cumulative and nondecreasing, $C_i(t+1) \geq C_i(t)$, then authority obtained from early contributions persists in full even after their contextual relevance has declined, producing institutional lock-in in which incumbency is self-reinforcing regardless of whether the incumbent's original domain of expertise still bears on present decisions. And the transfer-fidelity result compounds the concern further: even setting aside whether $C_i$ is canonical or decaying appropriately, the moment any resulting reward becomes a transferable token, authority stops tracking the responsibility that justified it at all, and instead tracks whoever currently holds the token.

A ledger without value replaces universal voting power with scoped standing. For a proposal $p$, the relevant set is not everyone possessing a balance but those with a defensible relation to the decision:
\[
\mathcal{S}(p) = \{ i \mid \mathrm{affected}(i,p) \lor \mathrm{responsible}(i,p) \lor \mathrm{competent}(i,p) \}.
\]
These predicates should remain distinct. Being affected grounds a right not to be ignored. Being responsible grounds accountability for implementation. Competence grounds evidential weight, not sovereignty. None should silently consume the others. Decision procedures can then vary by consequence: routine maintenance may use delegated responsibility; factual questions may use contestable expert judgment; allocation decisions may require consent from affected groups; constitutional changes may demand broad ratification and strong exit protections. The ledger supplies evidence about standing and history, but it does not compute a universal rank of persons.

\section{Contribution Without Permanent Reputation}

A lifetime contribution score can become a second currency: scarce, accumulable, unequally distributed, and convertible into access, and unlike money it may be impossible to give away or escape, producing hereditary institutions without hereditary titles, because early contributors indefinitely dominate the definition of later contribution. The alternative is contextual recognition with decay of authority but not decay of history. The record that an action occurred remains. Its relevance to a present decision must be argued again:
\[
\rho(c,p,t) = f(\text{domain fit}, \text{evidence}, \text{recency}, \text{affectedness}, \text{contestations}),
\]
where $\rho$ is not a stored personal score but a temporary assessment of how contribution $c$ bears on proposal $p$ at time $t$. Recency may serve as a useful proxy for current salience, but it is neither truth nor entitlement; older work remains reachable when similarity, dependency, or renewed relevance makes it operative again. This resembles tab completion more than a leaderboard: the system proposes plausible continuations from the current prefix and context, rather than declaring that the most historically prolific person owns the next symbol.

A worked scenario traces one contribution through the full architecture. A contributor, persona $P_{47}$, submits a patch addressing a scheduling bug: $\mathrm{Record}(e_{881})$. A CI run and a maintainer's manual review establish $\mathrm{Verify}(e_{881}, d_{\mathrm{scheduler}})$, a factual finding that would hold even if maintainers subsequently judged the fix undesirable for unrelated reasons.

Two maintainers attest that the fix is a genuine improvement, constituting
\[
\mathrm{Recognize}(P_{47}, e_{881}, d_{\mathrm{scheduler}}),
\]
scoped explicitly to that domain.

A governance process then authorizes payment, $\mathrm{Reward}(P_{47}, e_{881}, o_{19})$, recorded as an explicit obligation, not newly minted tokens. None of Record, Verify, Recognize, or Reward yet grants $P_{47}$ any say over unrelated decisions. Six months later, a proposal to change the project's licensing terms arises, a domain in which $P_{47}$ has no evidenced standing at all:
\[
\mathrm{Reward}(P_{47}, e_{881}, o_{19}) \not\Rightarrow \mathrm{Authorize}(P_{47}, p_{\mathrm{license}}).
\]
A separate proposal to change the scheduler's core allocation algorithm computes $P_{47}$'s standing as $\rho(c_{881}, p_{\mathrm{sched}}, t_{\mathrm{now}})$: domain fit is high, evidence is strong, recency has begun to decay, yielding meaningful evidential weight in the scheduler deliberation without granting unilateral control over the outcome. Contrast this with the original Ayian pipeline applied to the same event: the patch would generate a contribution score, the score would mint tokens, and $P_{47}$'s voting weight on the licensing proposal would be identical to their voting weight on the scheduler proposal, because the token does not know which domain it came from.

\section{Useful Work and Thermodynamic Honesty}

The phrase ``Proof of Contribution'' risks repeating the central mistake of proof-of-work systems: inventing measurable activity in order to manufacture trust. Useful computation is not whatever satisfies a protocol's puzzle. It is work that an independently existing human, scientific, ecological, or infrastructural need already required. Heat is the terminal by-product of work; a heater that produces heat directly discards the opportunity to perform useful computation or mechanical work along the way, and a cryptocurrency that invents computation solely to establish scarcity discards the opportunity to direct energy toward real unmet tasks, exactly the extraction structure Chapter~\ref{ch:ch19-ledgers-without-value} formalized as degeneracy: proof-of-work's native objective is orthogonal to any reduction of the work functional. A defensible system would schedule computation where heat is wanted, route tasks to appropriate climates and buildings, and count reduced duplication and recovered waste as part of its evidence.

But even genuinely useful work should not automatically mint governance. Solving a protein-folding problem, hosting winter computation, or maintaining a bridge may justify contractual payment and public acknowledgment. It does not create a natural claim to rule unrelated communities. The ledger can demonstrate that useful work occurred while refusing the metaphysical leap from energy expended to value owned. This is the six-act chain applied to a specific and increasingly common case: genuinely useful computational work is real Record and Verify, can readily earn Reward through an explicit obligation, and still does not by itself reach Authorize. Chapter~\ref{ch:ch11-constitutional-space-of-refusal} developed the constitutional layer this observation depends on, including the treatment of refusal itself as a first-class, non-punitive continuation option; Chapter~\ref{ch:ch22-sheaves-of-necessity} takes up the sheaf-theoretic generalization of the affordability and necessity questions this chapter's scoped-standing apparatus raises but does not itself resolve.
